\chapter{Mathematical Formulations and Biometric Metrics}

\section{Geometric Landmark Alignment via Affine Transformation}
Face recognition models degrade significantly in accuracy when exposed to in-plane head rotations, pitch variations, and scaling differences. To enforce geometric invariance, BioBytes subjects every detected facial region to a 2D Euclidean similarity transformation based on the five landmarks localized by YuNet.

Let the detected landmarks be denoted by $\mathbf{p}_i = (x_i, y_i)^T \in \mathbb{R}^2$ for $i \in \{1, \dots, 5\}$. A predefined canonical reference template defines standard landmark coordinates $\mathbf{q}_i \in \mathbb{R}^2$ on a $112 \times 112$ canonical canvas. The optimal transformation matrix $\mathbf{T} \in \mathbb{R}^{2 \times 3}$ minimizes the least-squares Euclidean error:
\begin{equation}
\min_{s, \mathbf{R}, \mathbf{t}} \sum_{i=1}^{5} \left\| \mathbf{q}_i - \left( s \mathbf{R} \mathbf{p}_i + \mathbf{t} \right) \right\|_2^2
\end{equation}
where $s \in \mathbb{R}^+$ represents the isotropic scale factor, $\mathbf{R} \in \text{SO}(2)$ is the 2D rotation matrix:
\begin{equation}
\mathbf{R}(\theta) = \begin{bmatrix} \cos\theta & -\sin\theta \\ \sin\theta & \cos\theta \end{bmatrix}
\end{equation}
and $\mathbf{t} = [t_x, t_y]^T \in \mathbb{R}^2$ is the translation vector.

The forward warp maps every target coordinate $(x', y')$ in the normalized crop back to the source frame:
\begin{equation}
\begin{bmatrix} x' \\ y' \\ 1 \end{bmatrix} = \begin{bmatrix} s\cos\theta & -s\sin\theta & t_x \\ s\sin\theta & s\cos\theta & t_y \\ 0 & 0 & 1 \end{bmatrix} \begin{bmatrix} x \\ y \\ 1 \end{bmatrix}
\end{equation}
Bilinear pixel interpolation is performed across the aligned bounding box, producing an aligned canonical image tensor $I_{\text{crop}} \in \mathbb{R}^{112 \times 112 \times 3}$.

\section{128-Dimensional Hypersphere Projection via SFace}
The canonical face crop $I_{\text{crop}}$ is ingested by the deep convolutional feature backbone $\mathcal{F}_{\text{SFace}}$:
\begin{equation}
\mathbf{z} = \mathcal{F}_{\text{SFace}}(I_{\text{crop}}) \in \mathbb{R}^{128}
\end{equation}
To eliminate sensitivity to absolute image brightness and contrast scaling, the raw feature vector $\mathbf{z}$ undergoes rigorous $\mathcal{L}_2$-norm projection onto a 128-dimensional unit hypersphere:
\begin{equation}
\mathbf{v} = \frac{\mathbf{z}}{\|\mathbf{z}\|_2} = \frac{\mathbf{z}}{\sqrt{\sum_{k=1}^{128} z_k^2}}
\end{equation}
Consequently, every enrolled and candidate biometric identity is constrained to the manifold:
\begin{equation}
\mathbf{v} \in \mathbb{S}^{127} = \left\{ \mathbf{x} \in \mathbb{R}^{128} : \|\mathbf{x}\|_2 = 1.0 \right\}
\end{equation}

\section{Metric Space Decision Rules: Cosine Similarity Invariant}
Given an operational probe vector $\mathbf{u} \in \mathbb{S}^{127}$ captured live and an enrolled student reference vector $\mathbf{v} \in \mathbb{S}^{127}$, identity verification is formulated in spherical Hilbert space. 

\subsection{Cosine Similarity Formulation}
Because all vectors possess unit norm ($\|\mathbf{u}\|_2 = \|\mathbf{v}\|_2 = 1$), the cosine similarity simplifies directly to their algebraic inner product:
\begin{equation}
S_C(\mathbf{u}, \mathbf{v}) = \frac{\mathbf{u} \cdot \mathbf{v}}{\|\mathbf{u}\|_2 \|\mathbf{v}\|_2} = \sum_{k=1}^{128} u_k v_k = \cos(\theta_{\mathbf{u}, \mathbf{v}})
\end{equation}
where $\theta_{\mathbf{u}, \mathbf{v}}$ is the geodesic angular displacement between the two identity representations on the hyperspherical manifold.

\subsection{Duality with Squared Euclidean Distance}
The relationship between cosine similarity and Euclidean distance on the unit hypersphere is strictly monotonic:
\begin{equation}
\|\mathbf{u} - \mathbf{v}\|_2^2 = \|\mathbf{u}\|_2^2 + \|\mathbf{v}\|_2^2 - 2(\mathbf{u} \cdot \mathbf{v}) = 2 - 2 S_C(\mathbf{u}, \mathbf{v})
\end{equation}
Hence:
\begin{equation}
D_E(\mathbf{u}, \mathbf{v}) = \sqrt{2 - 2 S_C(\mathbf{u}, \mathbf{v})}
\end{equation}

\subsection{The Decision Boundary and Invariant}
The classification function evaluates whether a probe belongs to enrolled student $j$ across a gallery of $N$ students:
\begin{equation}
j^* = \arg\max_{j \in \{1, \dots, N\}} S_C(\mathbf{u}, \mathbf{v}_j)
\end{equation}
The deterministic decision rule is parameterized by threshold $\tau$:
\begin{equation}
\text{Decision}(\mathbf{u}) = \begin{cases} 
\text{Authenticated as Student } j^*, & \text{if } S_C(\mathbf{u}, \mathbf{v}_{j^*}) \geq \tau \\
\text{Reject (Unenrolled / Impostor)}, & \text{if } S_C(\mathbf{u}, \mathbf{v}_{j^*}) < \tau
\end{cases}
\end{equation}
In BioBytes, empirical validation across diverse demographic test partitions establishes the invariant threshold at:
\begin{equation}
\tau = 0.363
\end{equation}

\section{Biometric Error Trade-Off (DET)}
To establish statutory confidence, system performance is modeled using continuous score distributions for genuine comparisons $p(s | \mathcal{H}_{\text{genuine}})$ and impostor comparisons $p(s | \mathcal{H}_{\text{impostor}})$:
\begin{itemize}
    \item \textbf{False Acceptance Rate (FAR):} Probability of erroneously accepting an impostor proxy:
    \begin{equation}
    \text{FAR}(\tau) = \int_{\tau}^{1} p(s \mid \mathcal{H}_{\text{impostor}}) \, ds
    \end{equation}
    \item \textbf{False Rejection Rate (FRR):} Probability of erroneously rejecting a legitimate enrolled student:
    \begin{equation}
    \text{FRR}(\tau) = \int_{-1}^{\tau} p(s \mid \mathcal{H}_{\text{genuine}}) \, ds
    \end{equation}
\end{itemize}
By enforcing $\tau = 0.363$, BioBytes achieves $\text{FAR} \leq 0.01\%$, satisfying the institutional mandate of zero-proxy tolerance.
